\subsection{méthode décompositoire}

Lorsque $A$ n'est pas symétrique, on peut utiliser la décomposition de LU~:
$$A = L U$$
où $L$ est triangulaire inférieure à diagonale unitaire (remplie de 1) et $U$ est triangulaire supérieure.\\
On applique un aglorithme basé sur le pivot de Gauss pour calculer cette factorisation, qui a une complexité de $O(\frac{2n^3}{3})$. La résolution du système linéaire à partir de cette factorisation est de $O(n^2)$.
\begin{lstLNat}
    U = $A$
    L = $I_n$
    pour i = 1 à n:
        // On veut que la colonne i de U soit nulle à partir de la ligne i+1.
        pour j = i+1 à n:
            alpha = U[j][i] / U[i][i]
            L[j][i] = alpha
            pour k = i à n:
                U[j][k] = U[j][k] - alpha * U[i][k]

\end{lstLNat}

\begin{exemple}{}{factorisation LU}
    Tentons de calculer la factorisation LU de la matrice~:
    $$A = \begin{pmatrix}
        2 & -1 & 3\\
        -4 & -1 & -2\\
        6 & -9 & 16
    \end{pmatrix}$$
    
    On obtient~:
    $$L = \begin{pmatrix}
        1 & 0 & 0\\
        -2 & 1 & 0\\
        3 & -3 & 1
    \end{pmatrix}$$
\end{exemple}

\begin{remarque}{}{sur cet algorithme}
    L'idée est de trigonaliser la matrice. Quand on cherche à annuler le coefficient $U[j][i]$ à l'aide de la ligne $i$, on stocke dans $L[j][i]$ l'opposé du rapport permettant de faire cette annulation, c'est-à-dire $-\frac{U[j][i]}{U[i][i]}$. Ainsi, à la fin de l'algorithme, on a $A = L U$.
\end{remarque}

\begin{remarque}{}{complexité}
    La complexité de la factorisation de LU est de $O(\frac{2n^3}{3})$, et la résolution du système linéaire à partir de cette factorisation est de $O(n^2)$. On néglige dans ces calculs la complexité du calcul des divisions, qui est de $O(1)$.
\end{remarque}

\begin{remarque}{}{inverse d'une matrice}
    Pour résoudre l'équation $Ax=b$, inverser $A$ est bien trop puissant, car cela implique déjà une une factorisation LU, et deux résolutions triangulaires.
\end{remarque}

\begin{remarque}{}{inverse d'une matrice}   
    Si $A$ est creuse, alors $A^{-1}$ est pleine.
\end{remarque}


\subsection{Méthode itérative}

Voir rapport projet 2, partie 2 pour la théories


\begin{lstLNat}
    x = $x_0$ // initialisation
    r = b - Ax // calcul du résidu
    k = 1
    % rho0 = $\code{r}^\top \code{r}$
    Tant que ($\norme{\code{r}} \norme{\code{b}}$ > \varepsilon et k < kmax):
        rho = $\code{r}^\top\code{r}$
         si k > 1:
            gamma = 0
        sinon:
            gamma = rho / rho\_old
        p = r + gamma * p
        
        Ap = $\code{A} \times \code{p}$ 
        alpha = rho / ($\code{p}^\top \code{Ap}$) // calcul du pas optimal
        x = x + alpha * p
        r = r - alpha * Ap
        rho\_old = rho
        k = k + 1
\end{lstLNat}

\begin{remarque}{}{précondtionneur}
    Pour prendre en compte le préconditionneur $M$, on ne doit pas avoir à manipuler $M$ ou son inverse. On va la manipuler sous sa décomposition de Cholesky $M = LL^\top$. On va alors faire les changements de variables suivants. On a alors~:
    \begin{align*}
        &z = M^{-1}r \\
        \iff Mz = r \\
        \iff LL^\top z = r\\
        \iff \begin{cases*}
            Ly = r\\
            L^\top z = y
        \end{cases*}
    \end{align*}
    
\end{remarque}
